\begin{abstract}
The viscoelastic gel-point of a polymer network is commonly treated as a fixed property of chemistry and concentration. Here we demonstrate, through three-dimensional random walk simulations on site-percolation lattices, that the gel-point is instead controlled by the growth dynamics of the forming network and can be continuously programmed via a single nucleation-density parameter. Probe-particle diffusion is modelled by random walkers confined to the unoccupied (void) sublattice of an $L \times L \times L$ cubic lattice ($L = 500$); the gel-point $p_c'$ is identified by the Winter--Chambon criterion $\alpha = 0.5$, where $\langle r^2(t) \rangle \propto t^\alpha$. For random (Bernoulli) percolation networks we find $p_c' \approx 0.683 \approx 1 - p_c$, quantitatively identifying the rheological gel-point with the loss of spanning connectivity in the void phase. Replacing random growth with adjacency-templated cluster growth shifts the gel-point to $p_c' \approx 0.886$ ($\Delta p_c' = +0.20$) and qualitatively alters the transition: the critical region is 3.4-times broader and probe particles remain persistently sub-diffusive ($\alpha \approx 0.3$--$0.4$) well above $p_c'$, a signature absent in the random class. Interpolating between these limits via a nucleation-density parameter $\kappa \in [0,1]$ — which controls the fraction of newly occupied sites placed as independent nuclei versus grown from the cluster frontier — yields a monotonically tunable gel-point spanning $\Delta p_c' \approx 0.31$ ($p_c' \in [0.683, 0.993]$). Frequency-domain viscoelastic spectra $G'(\omega)$ and $G''(\omega)$ derived via the Generalised Stokes--Einstein Relation confirm the $\delta(\omega) = 45^{\circ}$ Winter--Chambon signature at $p_c'$ for both gelation classes. These results establish nucleation density as a structural design parameter for programming the gel-point in network-forming materials.
\end{abstract}
